% AUTO-GENERATED by scripts/exp_sync_kernel_rate_curve_center_slice_audit.py
\begin{proposition}[中心截面 $Q(u)$ 的自倒易性与实根空性]\label{prop:rate-center-Q-reciprocal-no-real}
沿用证书 $R(\tfrac12,u)=-(u-1)^6Q(u)/64$ 中的 $Q(u)\in\ZZ[u]$。
由于 $Q$ 的系数严格回文，立得自倒易恒等式 $Q(u)=u^{16}Q(1/u)$，因此零点必以互逆对出现（并与复共轭合并为四元组）。
进一步地，对该整系数 $Q$ 做 Sturm 计数，得到其在实轴上的零点数为
$$
N_\RR(Q)=\,0.
$$
又因 $Q(0)=16>0$，故 $Q(u)>0$ 对所有 $u\in\RR$ 成立。
\end{proposition}
\leanverified{paper\_rate\_center\_q\_reciprocal\_no\_real}

\begin{corollary}[中心唯一实交与六重性]\label{cor:rate-center-unique-real}
\leanverified{paper\_rate\_center\_unique\_real}
在实参数 $u\in\RR$ 上，
$$
R\!\left(\tfrac12,u\right)=0\ \Longleftrightarrow\ u=1,
$$
且 $u=1$ 的重数恰为 $6$。
\end{corollary}

\begin{corollary}[中心截面的最近复奇点距离]\label{cor:rate-center-Rstar}
定义
$$
R_\star:=\min\{|u-1|:\ Q(u)=0\}.
$$
则数值上
$$
\boxed{\ R_\star\approx 1.0507607822\ }
$$
并可取最近根为
$$
u_\star\approx -0.0492815273\pm 0.0557359652i.
$$
\end{corollary}

\begin{corollary}[单位根扭曲的收敛域阈值（$p=5$ 被排除）]\label{cor:rate-center-omega-threshold}
\leanverified{paper\_rate\_center\_omega\_threshold}
令 $\omega_p=e^{2\pi i/p}$。则 $|\omega_p-1|=2\sin(\pi/p)$。
由
$$
|\omega_5-1|\approx 1.1755705046>R_\star,\qquad |\omega_7-1|\approx 0.8677674782<R_\star,
$$
并利用 $2\sin(\pi/p)$ 随 $p$ 单调递减，可得对所有 $p\ge 7$ 均有 $|\omega_p-1|<R_\star$。
因此任何围绕 $u=1$ 的局部解析展开，在 $u=\omega_p$ 处对 $p\ge 7$ 处于严格收敛域内，而在 $p=5$ 上必超出该域。
\end{corollary}

\begin{remark}[中心截面的极简审计链]\label{rem:rate-center-min-audit}
仅用打印出的 $Q$ 系数即可同时审计“对偶对称在中心截面的一致性”与“中心聚合次数”：
\begin{enumerate}
  \item 检查系数回文（等价于 $Q(u)=u^{16}Q(1/u)$）；
  \item 检查 $Q(1)=8515584\neq 0$（从而 $u=1$ 恰为 $6$ 重根）；
  \item 检查第六导数尺度（由分解式直接推出）
  $$
  Q(1)=-\frac{64}{6!}\,\partial_u^{6}R\!\left(\tfrac12,u\right)\Big|_{u=1}.
  $$
\end{enumerate}
\end{remark}
